% !TEX program = xelatex
\documentclass[11pt,a4paper,fontset=macnew]{ctexart}
\usepackage[left=20mm,right=20mm,top=18mm,bottom=18mm,headsep=5mm]{geometry}
\usepackage{fontspec}
\setmainfont{Times New Roman}
\usepackage{booktabs,longtable,array,enumitem,fancyhdr,needspace}
\usepackage[hidelinks,bookmarksopen=true]{hyperref}
\hypersetup{pdftitle={英语二作文：5套核心合并版},pdfsubject={合并同类文体，减少背诵负担；按用户给定批改标准}}
\newcolumntype{L}[1]{>{\raggedright\arraybackslash}p{#1}}
\setlength{\parindent}{0pt}
\setlength{\parskip}{4pt}
\setlength{\tabcolsep}{4pt}
\setlength{\emergencystretch}{2em}
\renewcommand{\arraystretch}{1.18}
\linespread{1.12}
\ctexset{section/format={\Large\bfseries},section/beforeskip=0pt,section/afterskip=8pt}
\setlist[itemize]{leftmargin=1.25em,itemsep=3pt,topsep=3pt}
\pagestyle{fancy}\fancyhf{}
\fancyhead[L]{\small 英语二作文 · 5套核心合并版}
\fancyhead[R]{\small\leftmark}
\fancyfoot[C]{\small\thepage}
\setlength{\headheight}{15pt}
\renewcommand{\headrulewidth}{0.3pt}
\newenvironment{english}{\begingroup\fontsize{11}{15}\selectfont\setlength{\parskip}{4pt}}{\par\endgroup}
\newenvironment{example}{\begingroup\fontsize{10.5}{14}\selectfont\setlength{\parskip}{4pt}}{\par\endgroup}
\begin{document}
\section*{英语二作文：5套核心合并版}
\addcontentsline{toc}{section}{背诵范围与合并关系}
\markboth{先明确背什么}{}
{\large\bfseries 主背第2—6页的5块框架，示例不用另背。}\par
\vspace{4pt}
上一版按36个类型分别写，重复的句子太多。本版把相似正文直接合并：开头只表达本题目的，中间只保留共用逻辑，格式与少数专用句另列查表。

\begin{longtable}{L{34mm}L{107mm}L{12mm}}\toprule
\textbf{只背这5块} & \textbf{合并后的用途} & \textbf{页码}\\\midrule\endhead
M1 活动安排 & 邀请、介绍、欢迎、通知、招募、计划、活动答复、倡议 & 2\\
M2 建议与理由 & 建议、推荐人或物、申请、求职、观点答复 & 3\\
M3 问题与补救 & 道歉、投诉、求助、辞职、婉拒 & 4\\
M4 情感交流 & 感谢、祝贺、表达感受、分享后追问 & 5\\
M5 图表作文 & 静态／动态共用后两段；多选、表格、双图只调整数据段 & 6\\
\bottomrule\end{longtable}
\par\vspace{5pt}\noindent\textbf{背诵量到底减少多少}\par\nopagebreak
五块框架共约196个固定英文词，已包含静态、动态两种开头；不含中文填空、称呼、署名和功能卡。第7页的功能卡是短句替换，不再扩成整篇。完整示例用来理解填什么，不增加一套背诵任务。

\par\vspace{5pt}\noindent\textbf{其余类型怎么处理}\par\nopagebreak
咨询信只用问句结构；寻物与招领共用“物品—特征—联系”；纪要、备忘录保留各自格式，见第8页。图画、文字与复合材料只留备用开头，见第9页。原来36类都有去向，不再各写一篇。

\par\vspace{5pt}\noindent\textbf{同一套怎么变成不同文章}\par\nopagebreak
邀请和介绍用M1同一个主体，只改变写作目的；道歉和投诉用M3同一个补救段，只改变责任与诉求；静态图和动态图用M5同一个解释段和结尾，只改变数据段。遇到“介绍＋邀请”“感受＋追问”时，把两个任务都写出来。

\par\vspace{5pt}\noindent\textbf{记住这三条就开始背}\par\nopagebreak
\begin{itemize}\item \textbf{先背框架，再换内容：}【】全部换成英文。动作通常用动词原形；doing提示处填动名词；“完整分句”要有主语和谓语。\item \textbf{正文仍按题目长度：}小作文常见约100词，大作文约150词。固定词少，是给具体信息留空间；填完后要核对总词数。\item \textbf{不能合并掉任务：}招募保留职责和要求，申请保留能力证据，道歉保留补救，图表保留准确数据。\end{itemize}
{\small 判断标准仍采用《考研英语批改指令》英语二部分：任务完整、语言准确、结构清楚、格式得体；大作文另核对数据和评论。}\par
\clearpage\section*{M1 活动安排｜同类只背这一块}\addcontentsline{toc}{section}{M1 活动安排｜同类只背这一块}\markboth{核心背诵}{}
{\small\textbf{合并：}邀请、介绍、欢迎、活动通知、招募、计划、活动类答复、倡议\par\textbf{思路：}为什么办 → 何时何地 → 做什么 → 读者怎么参与}\par
\par\vspace{5pt}\noindent\textbf{要背：固定框架}\par\nopagebreak
\begin{english}
The \textbf{【活动】} will take place \textbf{【时间】} at \textbf{【地点】}. Its aim is to \textbf{【目的，动词原形】}. I hope you can \textbf{【邀请或希望对方做的事，动词原形】}.

We plan to start with \textbf{【环节一】} and then move on to \textbf{【环节二】}. During the event, participants can \textbf{【动作】}, which will help them \textbf{【具体收获】}. Please \textbf{【准备或报名动作】} before \textbf{【时间／事件】}.

Let me know if you need any further details.

\end{english}
\par\vspace{5pt}\noindent\textbf{只看懂：填好后的同框架示例}\par\nopagebreak
{\small 例：邀请 Jack 参加校园旧物修理活动；说明安排和准备事项。\quad 正文99词。}\par
\begin{example}
Dear Jack,\par
The repair workshop will take place on Friday afternoon at the student centre. Its aim is to help students reuse old items. I hope you can join us and share your ideas.

We plan to start with a short demonstration and then move on to guided practice. During the event, participants can learn basic repair skills, which will help them make better use of what they own. Please bring a small item that needs repair before the workshop begins.

Let me know if you need any further details. I can reserve a place for you if you are interested.

Best wishes,\\Li Ming\par
\end{example}
\par\vspace{5pt}\noindent\textbf{换文体时，只改这些地方}\par\nopagebreak
{\small\begin{itemize}
\item 招募：第二段补职责和要求，见第7页，不另背通知正文。
\item 计划＋征求意见：把 will／plan to 改成 could，把结尾换为 What do you think of this plan?
\item 通知：去掉 I hope you can 和私人回信语，换标题、发布单位及报名方式。
\item 倡议：把活动目的换成行动目标；不必硬填时间地点，直接写两项可执行行动。
\end{itemize}}
\clearpage\section*{M2 建议与理由｜同类只背这一块}\addcontentsline{toc}{section}{M2 建议与理由｜同类只背这一块}\markboth{核心背诵}{}
{\small\textbf{合并：}建议、推荐人或物、申请、求职、观点类答复\par\textbf{思路：}提出选择 → 理由一＋具体作用 → 理由二／证据 → 下一步}\par
\par\vspace{5pt}\noindent\textbf{要背：固定框架}\par\nopagebreak
\begin{english}
To \textbf{【目标动作，动词原形】}, I suggest \textbf{【做法，动名词】}. There are two practical reasons for this choice.

First, \textbf{【做法或对象】} offers \textbf{【优势一】}. This can help \textbf{【受益者】} \textbf{【具体效果，动词原形】} without \textbf{【不必要的代价，动名词】}. Second, \textbf{【另一特点，完整分句】}. For example, \textbf{【贴题细节，完整分句】}. These features make it suitable for \textbf{【具体需要】}.

If you decide to try it, start by \textbf{【第一步，动名词】}.

\end{english}
\par\vspace{5pt}\noindent\textbf{只看懂：填好后的同框架示例}\par\nopagebreak
{\small 例：朋友 Anna 想提高口语，建议她参加小型讨论组，并说明理由。\quad 正文103词。}\par
\begin{example}
Dear Anna,\par
To improve your spoken English, I suggest joining a small discussion group. There are two practical reasons for this choice.

First, a small group offers regular speaking practice. This can help you express your ideas more confidently without paying for expensive lessons. Second, members can learn from one another. For example, when you cannot explain a point clearly, someone else may suggest a simpler expression. These features make it suitable for learners who need both practice and feedback.

If you decide to try it, start by inviting two classmates to discuss a familiar topic after class. I would be glad to join you.

Best wishes,\\Li Ming\par
\end{example}
\par\vspace{5pt}\noindent\textbf{换文体时，只改这些地方}\par\nopagebreak
{\small\begin{itemize}
\item 推荐：只改首句为 For【需求】, I recommend【人或物】；理由必须与对方需求匹配。
\item 直接给建议时，也可把 First 后改成一条行动，再保留 This can help... 解释好处。
\item 申请／求职：用第7页的申请开头；第二段写自己的技能＋具体经历，结尾请求面谈。
\item 推荐他人：补你与此人的关系，给真实情境中的能力证据，不只堆性格形容词。
\end{itemize}}
\clearpage\section*{M3 问题与补救｜同类只背这一块}\addcontentsline{toc}{section}{M3 问题与补救｜同类只背这一块}\markboth{核心背诵}{}
{\small\textbf{合并：}道歉、投诉、求助、辞职、婉拒\par\textbf{思路：}说明目的 → 问题与影响 → 解决或交接 → 确认}\par
\par\vspace{5pt}\noindent\textbf{要背：固定框架}\par\nopagebreak
\begin{english}
I am writing about \textbf{【事项】}. \textbf{【从第7页选择一句目的句】}

The main difficulty is that \textbf{【问题，完整分句】}. As a result, \textbf{【具体影响，完整分句】}. A workable solution would be to \textbf{【解决办法】}. In the meantime, \textbf{【临时安排，完整分句】}. This would allow \textbf{【受益者】} to \textbf{【预期结果】} without \textbf{【额外困难，动名词】}.

Please tell me whether this arrangement would work for you.

\end{english}
\par\vspace{5pt}\noindent\textbf{只看懂：填好后的同框架示例}\par\nopagebreak
{\small 例：你借 Alex 的书不慎损坏，说明影响并提出换书和临时借阅方案。\quad 正文103词。}\par
\begin{example}
Dear Alex,\par
I am writing about the book you lent me. I am sorry for damaging it yesterday.

The main difficulty is that several pages are now stained after I spilled water on them. As a result, some of the notes are difficult to read. A workable solution would be to buy you a new copy of the same edition. In the meantime, you can use my copy for your course. This would allow you to continue studying without waiting for the replacement to arrive.

Please tell me whether this arrangement would work for you. I will take better care of borrowed books in future.

Best wishes,\\Li Ming\par
\end{example}
\par\vspace{5pt}\noindent\textbf{换文体时，只改这些地方}\par\nopagebreak
{\small\begin{itemize}
\item 道歉要承担责任；投诉是客观说明问题并提出诉求，不能沿用“都是我的错”。
\item 辞职、婉拒必须在开头直接说清决定；补救段换成交接、替代帮助或改期。
\item 求助时写清自己遇到的困难、需要对方做的具体事情和期限。
\item The main difficulty... 不适合无实际困难的普通咨询；普通咨询直接用第7页问句。
\end{itemize}}
\clearpage\section*{M4 情感交流｜同类只背这一块}\addcontentsline{toc}{section}{M4 情感交流｜同类只背这一块}\markboth{核心背诵}{}
{\small\textbf{合并：}感谢、祝贺、表达感受、分享后追问、私人回复\par\textbf{思路：}明确情感 → 具体细节 → 为什么有意义 → 合适收尾}\par
\par\vspace{5pt}\noindent\textbf{要背：固定框架}\par\nopagebreak
\begin{english}
\textbf{【感谢／祝贺／回应消息的一句开头】}

What stayed with me was \textbf{【具体细节】}. It showed me how \textbf{【观察或感受，完整分句】}. I especially valued \textbf{【另一细节，名词或动名词】} because \textbf{【贴题理由，完整分句】}. This experience has encouraged me to \textbf{【积极行动，动词原形】}.

\textbf{【按任务选择：再次致谢／祝愿／具体追问】}

\end{english}
\par\vspace{5pt}\noindent\textbf{只看懂：填好后的同框架示例}\par\nopagebreak
{\small 例：Jack 分享了家庭旅行视频；表达具体感受，并询问想进一步了解的旅行细节。\quad 正文103词。}\par
\begin{example}
Dear Jack,\par
Your travel video made me smile, especially the scenes with your family.

What stayed with me was the moment when you learned to make noodles with a local cook. It showed me how food can bring people closer together. I especially valued the conversations with local families because they made the journey feel personal and warm. This experience has encouraged me to learn more about the places you visited.

How did you find the cooking workshop? I would also love to know which place your parents enjoyed most. Please share a few more stories when you have time. I would enjoy hearing them.

Best wishes,\\Li Ming\par
\end{example}
\par\vspace{5pt}\noindent\textbf{换文体时，只改这些地方}\par\nopagebreak
{\small\begin{itemize}
\item 感谢：细节写“对方具体帮了什么”，意义写“给我带来的变化”，最后再致谢即可。
\item 祝贺：细节写你看到的努力和成果；不适合写自身收获时删掉 This experience...。
\item 感受＋追问：一定写具体感受和具体问题；不能只用 Thank you／tell me more 凑齐任务。
\item 只有题目要求或情境合理时才写未来行动；这句可删，不必每种情感信都硬套。
\end{itemize}}
\clearpage\section*{M5 图表作文｜两种开头，共用后两段}\addcontentsline{toc}{section}{M5 图表作文｜两种开头，共用后两段}\markboth{核心背诵}{}
{\small\textbf{合并：}静态构成、多选柱图、群体表格、单指标趋势、多指标变化、波动、双图表}\par
\par\vspace{5pt}\noindent\textbf{要背：第一段二选一}\par\nopagebreak
\begin{example}\textbf{静态：类别／群体对比。}\par
The \textbf{【图表类型】} compares \textbf{【调查对象】} in terms of \textbf{【指标】}. \textbf{【A】} records the highest figure, at \textbf{【数值】}, while \textbf{【B】} stands at \textbf{【数值】}. \textbf{【另一关键对比，完整句】}

\textbf{动态：同一指标随时间变化。}\par
Between \textbf{【起年】} and \textbf{【止年】}, \textbf{【指标】} \textbf{【rose/fell】} from \textbf{【起值】} to \textbf{【终值】}. Over the same period, \textbf{【另一个指标或阶段的关键变化，完整分句】}. \textbf{【可由数据支持的总趋势，完整句】}

\end{example}
\par\vspace{5pt}\noindent\textbf{要背：第二、三段只保留这一套}\par\nopagebreak
\begin{example}
Two practical factors may help explain \textbf{【这种分布或变化，名词短语】}. First, \textbf{【原因一，完整分句】}. This can make it easier for \textbf{【群体】} to \textbf{【行为，动词原形】}. Second, \textbf{【原因二，完整分句】}. When \textbf{【具体情境，完整分句】}, people may \textbf{【选择或行动，动词原形】}. These are possible explanations, rather than causes established by the chart.

For \textbf{【有关主体】}, a useful next step would be to \textbf{【具体措施，动词原形】}. The aim should be to \textbf{【贴题目标】} while \textbf{【兼顾事项，动名词】}.

\end{example}
\par\vspace{5pt}\noindent\textbf{只看懂：动态示例，正文144词}\par\nopagebreak
{\small 新编数据：某市公共交通出行占比，2020年30\%、2022年36\%、2024年45\%；描述趋势并评论。}\par
\begin{example}
Between 2020 and 2024, the share of journeys made by public transport rose from 30\% to 45\%. Over the same period, the increase was larger after 2022, when the figure stood at 36\%. Overall, public transport gained fifteen percentage points.

Two practical factors may help explain this growth. First, better connections may have reduced the inconvenience of changing buses. This can make it easier for commuters to plan their journeys. Second, reliable arrival information may reduce uncertainty. When waiting times are clear, people may feel more comfortable leaving their cars at home. These are possible explanations, rather than causes established by the chart.

For city transport planners, a useful next step would be to improve connections and access. The aim should be to make daily travel easier while keeping fares affordable. Such improvements could also help people with different schedules benefit from the service.

\end{example}
{\small\textbf{只需调整：}表格和多选用静态开头；有转折时写关键升降，不能硬写steadily；双图各写一句核心信息。只有两个时间点时删掉“另一阶段”槽，改写整体增量。原因要能解释本图，不能一律填经济发展。}\par
\clearpage\section*{功能卡｜换文体，不新增整篇}\addcontentsline{toc}{section}{功能卡｜换文体，不新增整篇}\markboth{需要时查阅}{}
先选M1—M4，再从对应行选真正需要的一句。斜线表示选其一，不是把所有句子一起背、一起写。正式邮件和朋友邮件都可以直接进入主题。

{\small
\begin{longtable}{L{30mm}L{17mm}L{106mm}}\toprule
\textbf{任务} & \textbf{主干} & \textbf{只换这个短句／短语}\\\midrule\endhead
邀请／介绍／欢迎 & M1 & I hope you can join us. / The main arrangements are as follows.\\
招募要求 & M1 & Volunteers will \textbf{【职责】}. Applicants should \textbf{【条件】}.\\
推荐／申请 & M2 & I recommend \textbf{【对象】}. / I would like to apply for \textbf{【岗位】}.\\
道歉／投诉 & M3 & I am sorry for \textbf{【过失，动名词】}. / I would like \textbf{【解决诉求】}.\\
求助／辞职／婉拒 & M3 & Could you help me \textbf{【动作】}? / I have decided to leave \textbf{【岗位】}. / I cannot attend \textbf{【活动】}.\\
感谢／祝贺 & M4 & Thank you for \textbf{【帮助】}. / Congratulations on \textbf{【成就】}.\\
感受／追问 & M4 & Your \textbf{【分享内容】} made me feel \textbf{【感受】}. / Could you tell me \textbf{【疑问词＋陈述语序从句】}?\\
计划征求意见／申请面谈 & M1／M2 & What do you think of this plan? / I would welcome an interview.\\
\bottomrule\end{longtable}
}
\par\vspace{5pt}\noindent\textbf{三处不能只换一个单词}\par\nopagebreak
申请／求职：M2保留“理由＋证据”，把“推荐对象的优势”改成“我能带来的技能”，并加一句具体过去经历；结尾用面谈意愿。你是在证明自己适合岗位，不是在劝招聘者尝试一项活动。

道歉／投诉：M3相同的是问题、影响和解决办法。道歉须说明自己的过失；投诉须说明对方产品或服务的问题并提出诉求。不要把道歉的责任句带进投诉信。

计划／通知／招募：M1相同的是安排。共同计划用could并征求意见；通知直接发布事实；招募另外写“做什么工作、需要什么条件、如何报名”。

\par\vspace{5pt}\noindent\textbf{书信格式统一处理}\par\nopagebreak
称呼：Dear\textbf{【姓名／职务】},；朋友可用Best wishes,，正式信件可用Yours sincerely,；署名依题干，常用Li Ming。读者不认识你且身份影响任务时，加一句I am\textbf{【题目给定身份】}。不要凭空给自己加头衔。

{\small\textbf{词句只留够用的一组：}make it easier to...（使……更容易）；be suitable for...（适合……）；a workable solution（可行办法）；in the meantime（在此期间）；practical factors（实际因素）。都已出现在框架中，无须另背108条。}\par
\clearpage\section*{格式卡｜低频类型只记区别}\addcontentsline{toc}{section}{格式卡｜低频类型只记区别}\markboth{需要时查阅}{}
\par\vspace{5pt}\noindent\textbf{1. 通知／招募／倡议：主体借用M1}\par\nopagebreak
标题居中写NOTICE或具体倡议标题；正文直接写活动、行动或招募内容；结尾写报名办法和期限；落款按题目写组织或指定姓名。去掉Dear和Yours，不把通知写成私人邮件。

\par\vspace{5pt}\noindent\textbf{2. 咨询：同一问句结构，问两件具体的事}\par\nopagebreak
\begin{example}I am considering \textbf{【项目】}. Could you tell me \textbf{【疑问词＋陈述语序从句】}? I would also like to know whether \textbf{【完整分句】}. This would help me \textbf{【信息的用途，动词原形】}.

\end{example}
每个问题补一句相关背景或用途即可展开，按题目要求写到合适长度。例：Could you tell me when the course begins? 从句不用when does the course begin。

\par\vspace{5pt}\noindent\textbf{3. 寻物／招领：共用一条信息链}\par\nopagebreak
\begin{example}I \textbf{【lost/found】} \textbf{【物品】} \textbf{【时间和地点】}. It has \textbf{【可辨认的外观】}. Please \textbf{【联系或领取动作】} \textbf{【联系方式／地点】}.

\end{example}
寻物标题LOST，说明对自己的用途与联系办法；招领标题FOUND，说明保管、领取和认领核实办法。根据题目补足细节，不把这三句短框架直接当100词成文。

\par\vspace{5pt}\noindent\textbf{4. 纪要：会议信息＋决定＋谁何时完成}\par\nopagebreak
\begin{example}\textbf{Minutes of 【会议】}\par Date: 【日期】 \quad Place: 【地点】\par Present: 【与会者】\par
Members discussed \textbf{【议题】} and agreed to \textbf{【决定】}. \textbf{【负责人】} will \textbf{【任务】} by \textbf{【期限】}. Progress will be reviewed \textbf{【时间】}.

Recorded by: 【记录人】\par\end{example}
只记录会议实际讨论和决定的内容，按题目补充原因或分工。已发生的讨论用过去时，未来待办可用will；不能把纪要改成活动邀请。

\par\vspace{5pt}\noindent\textbf{5. 备忘录：只换信息栏，安排逻辑借用M1}\par\nopagebreak
\begin{example}\textbf{MEMO}\par To: 【对象】\quad From: 【发件人】\par Date: 【日期】\quad Subject: 【主题】\par\textbf{【变更或安排的事实】} Please \textbf{【要完成的动作】} by \textbf{【期限】}. Confirm completion when the task is finished.

\end{example}
正文补充原因、负责人及具体动作。纪要记“会上决定了什么”；备忘录告诉同事“现在需要做什么”。两者都不用再背一篇信。

\clearpage\section*{备用与核对｜不用增加背诵篇数}\addcontentsline{toc}{section}{备用与核对｜不用增加背诵篇数}\markboth{需要时查阅}{}
\par\vspace{5pt}\noindent\textbf{图画／文字／复合材料：保留入口，改掉图表专用句}\par\nopagebreak
{\small
\begin{longtable}{L{26mm}L{131mm}}\toprule
\textbf{材料} & \textbf{备用开头}\\\midrule\endhead
图画 & The picture shows \textbf{【人物】} \textbf{【动作，动名词】}. The contrast highlights \textbf{【主题】}.\\
文字材料 & The author argues that \textbf{【原观点】}. I agree with \textbf{【合理点】}, but \textbf{【自己的限制或不同意见】}.\\
图画＋图表 & The drawing shows \textbf{【场景】}, while the chart reports \textbf{【数据主题】}. \textbf{【核心数据，完整句】}\\
\bottomrule\end{longtable}
}
图画只在存在对比时用contrast。后面沿用“观点—理由—具体作用—例子—结论”的顺序；M5中的First、This can...、Second可以复用，但必须删掉causes established by the chart等图表专用句。文字题要回应是否同意，不能硬改成原因分析；复合材料要两部分都描述。

\par\vspace{5pt}\noindent\textbf{图表只保留六项核对}\par\nopagebreak
{\small
\begin{longtable}{L{32mm}L{125mm}}\toprule
\textbf{看什么} & \textbf{怎么避免写错}\\\midrule\endhead
对象、年份、单位 & 人次用visits，不是visitors；过去数据用过去时。\\
静态／动态 & 不同年龄组之间的差异，不等于同一群体随时间变化。\\
百分比／百分点 & 30\%到45\%：增加15个百分点；相对增长50\%。\\
最大／多数 & 45\%可为最大项；majority通常要求超过50\%。\\
多选／双图 & 多选不相加成总人数；双图并存不直接证明因果。\\
任务与词数 & 小作文约100词、大作文约150词，最终服从题干。\\
\bottomrule\end{longtable}
}
\par\vspace{5pt}\noindent\textbf{合并后的质量与文字重合检查}\par\nopagebreak
仍按你提供的《考研英语批改指令》第5—6页核对：不为难词牺牲准确性，不漏任务，不编图表数据，原因推测与可见事实分开。示例情境和数据为新编，不保证套上框架就有固定分数。

新版固定框架、示例和功能卡已分别对照原40份PDF的已提取文本检查。连续4词覆盖率最高17.65\%，低于50\%；这是字面复用指标，不是语义查重保证。计算以新稿英文词数为分母，命中的词位只计一次；中文占位和格式栏不计，扫描识别仍可能漏字。详细结果与源码保存在“源文件／合并减负版”。

{\small\textbf{来源取舍：}沿用上一版对南山任务拆解、各讲义功能句和英语二真题任务的整理；将重复表达归为5块，不再附36篇。完整原文件索引见上一版附录。原36类的逐项合并去向保存在本版content.json；408大纲仍不作为英语依据。}\par
\end{document}
